% =====================================================================
%  Harald Høffding, DANSKE FILOSOFER.  Kjøbenhavn og Kristiania:
%  Gyldendalske Boghandel, Nordisk Forlag, 1909 (MDCCCCIX).  206 pp.
%  TRANSCRIPTION (Danish) of the 1909 print, complete.
%
%  SOURCE: the Royal Library's (KB) scan, ~/bibliotek/Høffding, Harald/
%  danske-filosoffer.pdf (SCANS.tsv).  PAGE MAP (pagemap.py): printed page =
%  PDF page - 15; p. 1 = PDF 16 ... p. 206 = PDF 221.  Front matter
%  (unnumbered): half-title PDF 10, title PDF 12, contents PDF 14; these
%  carry "% --- front matter ---" comments and no \opage.  Not
%  transcribed: blank leaves, the publisher's pages after p. 206, folios,
%  running heads, sheet signatures („Harald Høffding: Danske Filosofer. 1“),
%  library stamps, pencil marks and underlinings in this copy.  Every page
%  turn is marked "% --- p. N ---" + \opage{N} at the first word of the
%  printed page; a word divided over a page turn ends in %.
%
%  CONVENTIONS: letterspacing (the print's emphasis, chiefly names) ->
%  \emph; italic -> \textit; printed dashes -> ---, number ranges -> --;
%  quotation marks „…“ as printed.  Chapters begin on a new page under a
%  numbered head in capitals („I. LUDVIG HOLBERG“, then „2.“-„17.“ as printed); the Indledning has no
%  head.  Footnotes at the page foot, marked *), **) on each page;
%  \footnote[k] gives the note's rank on its printed page and a comment
%  "% footnote *) on p. N" precedes each.  A note that runs over to the
%  next page is not split.
%
%  METHOD (the e-book method, 2026-10-02; TRANSCRIPTION-PLAYBOOK §00 point
%  0).  No one typed the text, except the Indledning (pp. 1-2), which the
%  SAGA e-book lacks: it is the scan's text layer read against the image.
%  For pp. 3-206 the base text -- words, punctuation, paragraphs, italics,
%  the footnote texts -- is the SAGA e-book (Danske_filosoffer.epub), which
%  sets the 115 footnotes as endnotes; ebook-method/danske-filosoffer/build.py
%  (tracked, repo root) aligned it
%  letter by letter with the scan's text layer, split per page by type
%  size into text and footnote lines, took every page turn from the layer
%  and every letterspaced word from the layer's spacing (names, verified
%  on a sample at the image; short words excluded).  Of 392 letter
%  disagreements 287 were settled by rule (the layer's glyph faults where
%  the e-book's word is attested elsewhere in the repository, letterspaced
%  capitals, specks, running heads, stamps, blocks the two streams place
%  differently); 105 letter, 91 punctuation and 20 paragraph disagreements
%  were decided at the image from crop sheets (decisions.tsv).  The print
%  corrects the e-book at: p. 14 Forgaaen; p. 17 Regering; p. 19 til;
%  p. 33 Fornuft; p. 45 Schein; p. 47 Oehlenschläger (twice); p. 61
%  nogetsomhelst; p. 78 sin; p. 79 til; p. 85 heldede; p. 97 forbigaaet;
%  p. 114 Bestaaen; p. 119 hans; p. 134 religionsfilosofiske; p. 137 für,
%  regnet; p. 178 dens; p. 193 externo; p. 90 n. den; and 27 places of
%  punctuation, chiefly commas the e-book adds („bøje, alt“, „sin,
%  frejdige“) or points it gives as commas.  PRINTED AS IS: p. 3 „at at“; p. 37
%  „Bamærkninger“ (corrected in pencil in this copy).
%
%  RESIDUAL RISKS: (1) a misprint that the e-book and the layer both
%  normalise is invisible to the method; (2) italics are the e-book's,
%  which the layer cannot witness; (3) letterspacing on words of fewer
%  than four letters is not recorded unless the word stands in a spaced
%  name; (4) p. 195, where the layer's line positions are unusable, had
%  no paragraph-indent check.
% =====================================================================
\documentclass[12pt,a4paper,openany]{book}
\usepackage[utf8]{inputenc}
\usepackage[T1]{fontenc}
\usepackage{textalpha}
\usepackage{libertinus}
\usepackage{libertinust1math}
\usepackage{graphicx}    % for \reflectbox (the old Danish ɔ: id-est mark)
\DeclareUnicodeCharacter{0254}{\reflectbox{c}}  % ɔ (U+0254) = reversed c
\usepackage[danish]{babel}
\usepackage[protrusion=true,expansion=true]{microtype}
\usepackage{setspace}
\usepackage[margin=1.2in]{geometry}
\usepackage{fancyhdr}
\setlength{\headheight}{28pt}
\addtolength{\topmargin}{-13.5pt}
\usepackage{hyperref}
\hypersetup{pdftitle={Danske Filosofer}, pdfauthor={Harald Høffding}, hidelinks}
\pagestyle{fancy}
\fancyhf{}
\fancyhead[LE,RO]{\thepage}
\fancyhead[RE]{\textit{Harald Høffding: Danske Filosofer}}
\fancyhead[LO]{\textit{1909}}
\setstretch{1.3}
\usepackage{marginnote}
\newcommand{\opage}[1]{\marginnote{\footnotesize\textit{[#1]}}}
% Footnotes as the print marks them, per page: *), **), ***), †).
\renewcommand{\thefootnote}{\ifcase\value{footnote}\or*)\or**)\or***)\or†)\fi}
\newcommand{\tocline}[3]{\noindent\makebox[2.2em][r]{#1}\hspace{0.6em}#2\dotfill\ #3\par}

\begin{document}

% --- front matter: half-title (PDF 10) ---
\begin{center}\vspace*{6em}DANSKE FILOSOFER\end{center}
\clearpage

% --- front matter: title (PDF 12) ---
\begin{center}
{\large HARALD HØFFDING}\\[6pt]
\rule{2cm}{0.4pt}\\[12pt]
{\Huge DANSKE}\\[6pt]
{\Huge FILOSOFER}\\[8em]
% (the publisher's device)
{\small GYLDENDALSKE BOGHANDEL}\\
{\small NORDISK FORLAG}\\[2pt]
\rule{0.6cm}{0.4pt}\\
{\scriptsize MDCCCCIX}
\end{center}
\clearpage

% --- front matter: contents (PDF 14) ---
\begin{center}{\large INDHOLD}\end{center}
\medskip
\tocline{}{Indledning}{1}
\tocline{1.}{Ludvig Holberg}{3}
\tocline{2.}{Sneedorf og Wolfianerne}{16}
\tocline{3.}{Danske Kantianere og deres Modstandere}{25}
\tocline{4.}{Henrik Steffens}{40}
\tocline{5.}{Hans Christian Ørsted}{49}
\tocline{6.}{Grundtvig og H. C. Ørsted}{57}
\tocline{7.}{Baggesen og Grundtvig}{66}
\tocline{8.}{Anders Sandøe Ørsted}{73}
\tocline{9.}{F. G. Howitz}{81}
\tocline{10.}{Niels Treschow}{89}
\tocline{11.}{F. C. Sibbern}{97}
\tocline{12.}{Poul Møller}{118}
\tocline{13.}{Heiberg og Martensen}{128}
\tocline{14.}{Søren Kierkegaard}{147}
\tocline{15.}{Frederik Dreier}{175}
\tocline{16.}{Rasmus Nielsen}{184}
\tocline{17.}{Hans Brøchner}{196}
\clearpage

% ===== INDLEDNING (pp. 1--2): not in the e-book; the scan's text layer, read at the image =====
% --- p. 1 ---
\setcounter{footnote}{0}%
\opage{1}ET lille Land som vort ligger aabent for Indflydelser udefra. De
aandelige Strømme fra de store Kulturlande gribe bestemmende ind i dets
Tankeliv, og dette kan ikke faa den faste Sammenhæng gennem Tiderne, som er
muligt, naar et Land kan frembringe en stor Mangfoldighed af tænkende Aander,
der aldrig lade Udviklingen stanse og ved de forskellige Typer, de frembyde,
kunne bryde Paavirkningerne udefra, saa at Sammenhængen i Folkets Tankeliv
bedre kan bevares.

Alligevel er der blevet tænkt alvorligt og selvstændigt herhjemme om de
Spørgsmaal, det er Filosofiens Kald at drøfte, --- om Videnskabens
Forudsætninger, om dens Resultaters Betydning for Livsanskuelsen, om
Grundlaget for etisk Vurdering, om den menneskelige Aands Natur og Love, om
Forholdet mellem Religion og Humanitet. De ydre Strømninger ere blevne
tilegnede paa ejendommelig Maade, og vor Nationalkarakter har lagt sig for
Dagen i den Maade, paa hvilken Problemerne ere blevne behandlede. Ved
Behandlingen af saadanne Spørgsmaal som de anførte, der ligge paa Grænserne
mellem den strenge Videnskab og Livet, maa Tænkernes Personligheder
nødvendigvis medvirke i højere Grad end ved andre Undersøgelser. Og hvor
Personligheden virker efter sin hele Ejendommelighed, gøre ogsaa de nationale
Egenskaber sig gældende. De ere jo ikke Lag eller Stykker, der udefra ere
føjede til de Enkeltes andre Egenskaber, men bestemme fra først af disse og
klinge med, hvergang Sindets Strenge tone.

At dansk Tænkning har sin Ejendommelighed, kan ses ved Sammenligning med
svensk Tænkning. Sveriges Aandsliv
% --- p. 2 ---
\setcounter{footnote}{0}%
\opage{2}har modtaget de samme Strømninger fra Udlandet som Danmarks, men de ere paa
Tankens Omraade blevne tilegnede og udviklede paa anden Maade og i anden
Retning end hos os. Spekulation og Mystik have spillet en langt større Rolle i
Sverige end i Danmark, og i Sammenhæng dermed har en Trang til systematisk
Afslutning gjort sig gældende, der aldrig har fundet fast Jordbund hos os. Den
danske Tænkning har væsentligt været baaret af psykologisk og etisk Interesse,
og der har atter og atter gjort sig en besindig Kritik gældende overfor
Systemdannelserne.%
% footnote *) on p. 2
\footnote[1]{Se min Afhandling om svensk og dansk Filosofi i Samlingen
\textit{Norden} (1902), saavel som min Afhandling \textit{Filosofien i Sverige}
(Det Letterstedtske Tidsskrift 1879). --- Om de filosofiske Bevægelser i
Udlandet, der faa Indflydelse paa dansk Filosofi, vil man finde nærmere
Oplysning i mit Værk \textit{Den nyere Filosofis Historie} ved Hjælp af
Registret. --- Biografier giver jeg som Regel ikke, især fordi jeg kan henvise
til de fortrinlige Artikler af min Kollega Professor \emph{Kroman} i
„Biografisk Leksikon“.} Hos vore ypperste tænkende Aander (Holberg,
Ørstederne, Sibbern, Kierkegaard) gør denne Ejendommelighed sig gældende, og
den afføder Tanker og Værker, der vel ikke have grebet ind i den store
internationale Diskussion om Erkendelsens og Livets Spørgsmaal, men som dog
have deres Værdi --- og det ikke blot for os. En Karakteristik af danske
Filosofer vil kunne lære os Mangt og Meget til Orientering overfor hine
Spørgsmaal, og ofte vil det vise sig, at vore Tænkere meget vel vilde have
kunnet gøre sig gældende paa en større Skueplads, om Forholdene havde anvist
dem en saadan. Tankens Himmel er, ligesom den fysiske Himmel, tilgængelig for
Alle, ikke blot for dem, der fødtes i et stort Folk. Denne Bog vil derfor kunne
kaste Lys over de filosofiske Problemer, ikke blot tjene til Karakteristik af
dansk Aandsliv saaledes som det ytrer sig paa Tænkningens Omraade.

\begin{center}\rule{2cm}{0.4pt}\end{center}
\clearpage

% ===== TEXT (pp. 3--206) =====
%%BODY%%

\end{document}
